\chapter{Metamodel kontekstowy W20-Binomial}

\section{Motywacja budowy metamodelu}

\subsection{Rola modelu drugiego poziomu}

\subsubsection{Ograniczenia pojedynczych rankingów}

W poprzednim rozdziale pokazano, że systemy rankingowe zawodników, szczególnie Player Glicko-2, dostarczają silnego sygnału predykcyjnego. Pojedynczy ranking nie wykorzystuje jednak pełnej informacji dostępnej przed meczem. Nie łączy jednocześnie wielu rodzin ratingów, miar niepewności, historycznej formy drużyn oraz formatu serii Bo1/Bo3/Bo5. Z tego powodu kolejnym etapem pracy jest metamodel, którego zadaniem jest połączenie tych informacji w jedną predykcję probabilistyczną.

\subsubsection{Stacking jako uzasadnienie konstrukcji modelu}

Metamodel jest w tej pracy rozumiany jako model drugiego poziomu. Jego wejściami są predykcje i parametry modeli rankingowych oraz cechy historyczne obliczane wyłącznie na podstawie spotkań rozegranych przed przewidywanym meczem. Takie podejście jest zgodne z ideą stacked generalization, w której predykcje modeli bazowych są wykorzystywane jako cechy dla modelu wyższego poziomu \parencite{wolpert1992,breiman1996stacked}.

\subsection{Definicja wariantu W20-Binomial}

Finalnym obiektem modelowania w tym rozdziale jest \textbf{W20-Binomial}: kompletny zestaw cech obejmujący predykcje rankingowe, niepewności rankingów, pełny historyczny kontekst drużynowy W20 oraz sygnały rankingowe skorygowane dwumianowo względem formatu serii. Rodzina estymatora nie została założona z góry. Hiperparametry kilku rodzin modeli dobrano za pomocą Optuny, a następnie dostrojone warianty porównano w procedurze walk-forward i miesięcznym bootstrapie blokowym. Ostatecznie najstabilniejszym estymatorem okazała się regresja logistyczna ElasticNet.

\section{Zestaw cech W20-Binomial}

Wszystkie cechy dla meczu $t$ były obliczane wyłącznie z informacji dostępnych przed tym meczem. Model nie korzystał z żadnych statystyk z rozgrywanego spotkania. Zestaw W20-Binomial składa się z czterech grup cech.

\subsection{Predykcje rankingowe zawodników}

Pierwszą grupę stanowią \textbf{predykcje rankingowe zawodników}, czyli prawdopodobieństwa zwycięstwa Team 1 wygenerowane przez sześć rodzin rankingowych:

\begin{itemize}
    \item \texttt{player\_elo} -- Elo,
    \item \texttt{player\_gl} -- Glicko-2,
    \item \texttt{player\_ts} -- TrueSkill,
    \item \texttt{player\_os} -- OpenSkill,
    \item \texttt{player\_pl} -- Plackett--Luce,
    \item \texttt{player\_tm} -- Thurstone--Mosteller.
\end{itemize}

\subsection{Niepewność ocen i rozkład siły składu}

Drugą grupę stanowią \textbf{cechy niepewności i rozkładu siły w składzie}. Obejmują one między innymi minimalne ratingi zawodników w Elo, maksymalne ratingi w Glicko-2, średnie odchylenie ratingu Glicko-2 oraz średnie niepewności $\sigma$ dla rodzin TrueSkill, OpenSkill, Plackett--Luce i Thurstone--Mosteller. Celem tych cech jest opisanie nie tylko średniej siły drużyny, lecz także stabilności i pewności tej oceny.

\subsection{Historyczny kontekst drużynowy W20}

Trzecią grupę stanowi \textbf{pełny historyczny kontekst W20}. Dla obu stron meczu agregowano ostatnie 20 wcześniejszych obserwacji danej drużyny. Kontekst obejmuje:

\begin{itemize}
    \item kroczący współczynnik zwycięstw,
    \item średnią liczbę zabójstw i śmierci,
    \item historyczne GD@15,
    \item DPM i VSPM,
    \item liczbę zniszczonych wież,
    \item liczbę Baronów Nashora,
    \item końcową ilość złota drużyny,
    \item czas trwania gry.
\end{itemize}

W dalszej części pracy przez \textbf{model kontekstowy} rozumiany jest model wykorzystujący właśnie ten pełny kontekst historyczny. Nie rozdzielano finalnej narracji na wariant podstawowy i pełny; węższe zestawy cech traktowano wyłącznie jako ablacje diagnostyczne.

\subsection{Cechy dwumianowe dla formatu serii}

Czwartą grupę stanowią \textbf{cechy dwumianowe} (ang. \textit{binomial features}). Format serii Bo1/Bo3/Bo5 nie został podany modelowi wyłącznie jako surowa liczba. Zamiast tego prawdopodobieństwa rankingowe przekształcono do przybliżonego prawdopodobieństwa wygrania całej serii. Dla serii do $n$ map zastosowano wzór:

\[
    P(\text{win series}) = \sum_{k=\lfloor n/2 \rfloor + 1}^{n}
    \binom{n}{k} p^k (1-p)^{n-k},
\]

gdzie $p$ oznacza rankingowe prawdopodobieństwo wygrania pojedynczej mapy lub pojedynczego meczu traktowane jako sygnał bazowy. W praktyce utworzono po jednej cesze dwumianowej dla każdej rodziny rankingowej. Zabieg ten jest istotny, ponieważ dłuższa seria nie tylko informuje o typie meczu, lecz także zmienia wariancję wyniku: silniejsza drużyna powinna mieć większą szansę wygrania Bo3 lub Bo5 niż pojedynczej mapy.

\begin{table}[H]
    \centering
    \small
    \caption{Skład finalnego zestawu cech W20-Binomial}
    \label{tab:metamodel_feature_sets}
    \begin{tabular}{p{3.4cm}p{7.3cm}p{3.3cm}}
        \toprule
        Grupa cech & Skład & Rola w modelu \\
        \midrule
        Predykcje rankingowe & Elo, Glicko-2, TrueSkill, OpenSkill, Plackett--Luce, Thurstone--Mosteller & Bazowy opis relatywnej siły składów \\
        Niepewność rankingów & Minimalne/maksymalne ratingi, RD Glicko-2, średnie $\sigma$ rodzin probabilistycznych & Informacja o stabilności i pewności oceny \\
        Kontekst W20 & Win rate, kills, deaths, GD@15, DPM, VSPM, towers, nashors, gold, duration & Forma, styl, ekonomia, obiekty i tempo gry \\
        Korekta dwumianowa & Sześć prawdopodobieństw rankingowych przekształconych wzorem dwumianowym dla Bo1/Bo3/Bo5 & Strukturalne uwzględnienie formatu serii \\
        \bottomrule
    \end{tabular}
\end{table}

\begin{figure}[H]
    \centering
    \resizebox{0.92\textwidth}{!}{%
    \begin{tikzpicture}[
        block/.style={draw, rounded corners, align=center, minimum width=4.0cm, minimum height=1.0cm, fill=blue!6},
        feature/.style={draw, rounded corners, align=center, minimum width=4.0cm, minimum height=1.0cm, fill=green!7},
        process/.style={draw, rounded corners, align=center, minimum width=5.2cm, minimum height=1.0cm, fill=orange!10},
        output/.style={draw, rounded corners, align=center, minimum width=4.4cm, minimum height=1.0cm, fill=purple!8},
        arrow/.style={-{Latex[length=2.5mm]}, thick, shorten >=2pt, shorten <=2pt}
    ]
        \node[block] (golgg) at (0,0) {Dane GOL.GG\\mecze, składy, statystyki};

        \node[feature] (rankings) at (-3.2,-2.2) {Systemy rankingowe\\Elo, Glicko-2, TS, OS, PL, TM};
        \node[feature] (context) at (3.2,-2.2) {Pełny kontekst W20\\forma, ekonomia, obiekty, tempo};

        \node[process] (binom) at (-3.2,-4.4) {Transformacja dwumianowa\\zależnie od długości serii};
        \node[process] (features) at (0,-6.6) {Wektor cech\\W20-Binomial};
        \node[process] (lr) at (0,-8.8) {Regresja logistyczna ElasticNet};
        \node[process] (sym) at (0,-11.0) {Symetryzacja kolejności drużyn\\$p_{sym}=\frac{p_{orig}+(1-p_{swap})}{2}$};
        \node[process] (cal) at (0,-13.2) {Kalibracja Platta\\expanding-window};
        \node[output] (prob) at (0,-15.4) {Końcowa predykcja\\$P(\mathrm{Team\ 1\ wins})$};

        \draw[arrow] (golgg.south) to[out=-135,in=90] (rankings.north);
        \draw[arrow] (golgg.south) to[out=-45,in=90] (context.north);

        \draw[arrow] (rankings.south) -- (binom.north);

        \draw[arrow] (context.south) to[out=-90,in=35] (features.east);
        \draw[arrow] (binom.south) to[out=-90,in=145] (features.north);

        \draw[arrow] (features.south) -- (lr.north);
        \draw[arrow] (lr.south) -- (sym.north);
        \draw[arrow] (sym.south) -- (cal.north);
        \draw[arrow] (cal.south) -- (prob.north);
    \end{tikzpicture}%
    }
    \caption{Architektura finalnego modelu W20-Binomial. Diagram pokazuje cztery główne grupy informacji wejściowej oraz dwa końcowe etapy postprocessingu: symetryzację predykcji względem kolejności drużyn i kalibrację Platta w schemacie expanding-window.}
    \label{fig:w20_binomial_architecture}
\end{figure}

\section{Współzależność sygnałów rankingowych}

\subsection{Korelacje pomiędzy rodzinami rankingów}

Przed trenowaniem modelu sprawdzono, czy sygnały rankingowe rzeczywiście wnoszą niezależną informację. Wiele systemów rankingowych opisuje tę samą ukrytą wielkość: relatywną siłę drużyny lub zawodników przed meczem. Z tego powodu można oczekiwać wysokich korelacji pomiędzy predykcjami Elo, Glicko-2, TrueSkill, OpenSkill, Plackett--Luce i Thurstone--Mosteller.

\begin{figure}[H]
    \centering
    \includegraphics[width=0.9\textwidth]{metamodel_correlation_vif/rating_signal_correlation_heatmap.png}
    \caption{Macierz korelacji głównych sygnałów rankingowych.}
    \label{fig:rating_signal_correlation}
\end{figure}

Macierz na rysunku~\ref{fig:rating_signal_correlation} pokazuje bardzo silną współzależność wielu ratingów. Nie oznacza to, że cechy rankingowe są bezużyteczne; przeciwnie, pokazuje, że zawierają podobny sygnał sportowy. Oznacza jednak, że model liniowy powinien być regularyzowany, a interpretacja pojedynczych współczynników musi być ostrożna.

\subsection{Analiza VIF i konsekwencje dla modelu liniowego}

Drugim krokiem była analiza VIF, czyli variance inflation factor. Wysoki VIF wskazuje na silną współliniowość i może prowadzić do niestabilnych współczynników w modelach liniowych.

\begin{table}[H]
    \centering
    \caption{Najwyższe wartości VIF dla sygnałów rankingowych}
    \label{tab:rating_signal_vif}
    \begin{tabular}{lr}
        \toprule
        Sygnał rankingowy & VIF \\
        \midrule
        Team Thurstone--Mosteller & 337{,}6 \\
        Team OpenSkill & 222{,}3 \\
        Team Glicko-2 & 143{,}4 \\
        Player OpenSkill & 132{,}6 \\
        Player Thurstone--Mosteller & 126{,}4 \\
        Player Plackett--Luce & 124{,}5 \\
        Team TrueSkill & 109{,}4 \\
        Team Plackett--Luce & 106{,}3 \\
        Player Elo & 59{,}7 \\
        Team Elo & 41{,}4 \\
        Player TrueSkill & 25{,}4 \\
        Player Glicko-2 & 13{,}2 \\
        \bottomrule
    \end{tabular}
\end{table}

Wyniki w tabeli~\ref{tab:rating_signal_vif} potwierdzają, że część sygnałów rankingowych jest silnie współliniowa. W pracy nie potraktowano tego jako powodu do ręcznego usunięcia całych rodzin ratingów, ponieważ podobne systemy mogą nadal wnosić niewielkie, ale użyteczne różnice w konkretnych okresach, ligach lub składach. Wniosek z analizy VIF jest zatem metodologiczny: finalny metamodel powinien wykorzystywać regularyzację oraz powinien być oceniany przede wszystkim przez jakość predykcji out-of-time, a nie przez bezpośrednią interpretację pojedynczych współczynników.

Z tego powodu w dalszych eksperymentach zastosowano regresję logistyczną ElasticNet oraz porównano ją z modelami nieliniowymi w procedurze walk-forward. Regularyzacja ElasticNet ogranicza wpływ współliniowości przez zmniejszanie wag podobnych cech i częściową selekcję sygnałów, natomiast bootstrap blokowy sprawdza, czy uzyskana przewaga nie wynika z jednej niestabilnej konfiguracji danych. Analiza VIF pełni więc rolę uzasadnienia wyboru regularyzowanego estymatora i ostrożnej interpretacji wag modelu, a nie osobnego kryterium eliminacji cech.

\section{Dobór długości okna historycznego}

Osobno sprawdzono długość okna historycznego dla pełnego kontekstu drużynowego. Testowano okna W3, W5, W10, W20 oraz W50. Krótsze okna szybciej reagują na zmianę formy, lecz są bardziej podatne na szum pojedynczych meczów. Dłuższe okna są stabilniejsze, ale mogą mieszać aktualną formę z nieaktualnymi składami, patchami lub stylami gry.

\begin{table}[H]
    \centering
    \caption{Wpływ długości okna historycznego na jakość metamodelu}
    \label{tab:metamodel_window_ablation}
    \begin{tabular}{lrrrr}
        \toprule
        Okno & AUC & LogLoss & Brier & ECE \\
        \midrule
        W3 & 0{,}7470 & 0{,}5928 & 0{,}2036 & 0{,}0131 \\
        W5 & 0{,}7476 & 0{,}5924 & 0{,}2035 & 0{,}0108 \\
        W10 & 0{,}7477 & 0{,}5925 & 0{,}2034 & 0{,}0088 \\
        W20 & 0{,}7492 & 0{,}5910 & 0{,}2028 & 0{,}0088 \\
        W50 & 0{,}7489 & 0{,}5915 & 0{,}2030 & 0{,}0124 \\
        \bottomrule
    \end{tabular}
\end{table}

Najniższy LogLoss w tej ablacji uzyskało okno W20. Różnica względem W50 nie jest duża, ale W20 lepiej odpowiada finalnej interpretacji modelu: ma reprezentować aktualną formę, styl i tempo drużyny, bez nadmiernego obciążenia starszymi meczami. Dlatego w dalszych eksperymentach jako główny wariant kontekstu przyjęto W20.

\section{Dobór rodziny estymatora po strojeniu hiperparametrów}

Ostatnim krokiem wyboru architektury było porównanie kilku rodzin modeli na tym samym kompletnym zestawie cech W20-Binomial. Hiperparametry dla każdej rodziny dobrano za pomocą Optuny na chronologicznym zbiorze walidacyjnym. Wyników samej walidacji Optuny nie traktowano jako finalnej tabeli rezultatów; Optuna pełniła funkcję strojenia hiperparametrów. Dopiero dostrojone warianty oceniono w schemacie walk-forward, który lepiej odpowiada rzeczywistemu użyciu modelu w czasie.

Porównano regresję logistyczną ElasticNet, ExtraTrees, LightGBM, HistGradientBoosting oraz małą sieć MLP. Wszystkie modele korzystały z tego samego zestawu cech W20-Binomial i były trenowane w schemacie expanding-window z aktualizacją co 1000 meczów.

\begin{table}[H]
    \centering
    \caption{Porównanie dostrojonych rodzin modeli W20-Binomial w procedurze walk-forward}
    \label{tab:w20_binomial_model_family_comparison}
    \begin{tabular}{lrrrr}
        \toprule
        Model & AUC & LogLoss & Brier & ECE \\
        \midrule
        LR ElasticNet & 0{,}7553 & 0{,}5858 & 0{,}2007 & 0{,}0130 \\
        ExtraTrees & 0{,}7505 & 0{,}5896 & 0{,}2024 & 0{,}0121 \\
        LightGBM & 0{,}7507 & 0{,}5897 & 0{,}2024 & 0{,}0110 \\
        HistGradientBoosting & 0{,}7502 & 0{,}5904 & 0{,}2026 & 0{,}0100 \\
        MLP & 0{,}7456 & 0{,}5952 & 0{,}2046 & 0{,}0201 \\
        \bottomrule
    \end{tabular}
\end{table}

Najniższy LogLoss i najwyższy AUC uzyskała regresja logistyczna ElasticNet. Modele nieliniowe miały miejscami lepsze ECE, lecz zgodnie z hierarchią metryk głównym kryterium wyboru pozostaje LogLoss, ponieważ bezpośrednio mierzy jakość prognozy probabilistycznej i silnie karze błędną pewność.

\begin{figure}[H]
    \centering
    \includegraphics[width=0.92\textwidth]{w20_binomial_all_models_bootstrap/all_models_bootstrap_delta_logloss.png}
    \caption{Miesięczny bootstrap blokowy różnic LogLoss dostrojonych modeli względem regresji logistycznej ElasticNet.}
    \label{fig:w20_binomial_all_models_bootstrap}
\end{figure}

Bootstrap blokowy przedstawiony na rysunku~\ref{fig:w20_binomial_all_models_bootstrap} potwierdza stabilność przewagi modelu liniowego. Wszystkie testowane warianty nieliniowe miały dodatnie przedziały ufności dla różnicy LogLoss względem regresji logistycznej ElasticNet. Różnice wynosiły odpowiednio około 0{,}0038 dla ExtraTrees, 0{,}0039 dla LightGBM, 0{,}0045 dla HistGradientBoosting oraz 0{,}0094 dla MLP. Oznacza to, że wybór regresji logistycznej nie wynika wyłącznie z pojedynczej tabeli metryk, lecz jest stabilny względem miesięcznego resamplingu.

\section{Ablacje finalnego estymatora: masking i sygnał Glicko-2}

Po wyborze regresji logistycznej ElasticNet jako finalnej rodziny estymatora przeprowadzono dodatkowe ablacje wyjaśniające, z których sygnałów korzysta model. Taka kolejność jest istotna: najpierw wybrano klasę modelu na kompletnym zestawie cech W20-Binomial, a dopiero następnie sprawdzono wrażliwość wybranego estymatora na maskowanie i usuwanie wybranych grup rankingowych. Szczególnie ważny był test zależności od rodziny Glicko-2, ponieważ w rozdziale rankingowym Player Glicko-2 był najlepszym pojedynczym sygnałem.

\begin{table}[H]
    \centering
    \caption{Masking i ablacja sygnału Player Glicko-2 w modelu kontekstowym}
    \label{tab:logistic_masking_ablation}
    \begin{tabular}{lrrrr}
        \toprule
        Wariant & AUC & LogLoss & ECE & $\Delta$ LogLoss \\
        \midrule
        Pełny model kontekstowy & 0{,}7555 & 0{,}5855 & 0{,}0118 & 0{,}0000 \\
        G2 + context & 0{,}7551 & 0{,}5861 & 0{,}0102 & 0{,}0005 \\
        Mask all 10\% & 0{,}7537 & 0{,}5868 & 0{,}0085 & 0{,}0012 \\
        Mask G2 10\% & 0{,}7527 & 0{,}5875 & 0{,}0089 & 0{,}0020 \\
        Mask G2 30\% & 0{,}7516 & 0{,}5884 & 0{,}0110 & 0{,}0028 \\
        No player\_gl & 0{,}7501 & 0{,}5898 & 0{,}0130 & 0{,}0042 \\
        No G2 family & 0{,}7470 & 0{,}5925 & 0{,}0119 & 0{,}0070 \\
        \bottomrule
    \end{tabular}
\end{table}

Eksperyment potwierdza, że Player Glicko-2 jest najważniejszym pojedynczym sygnałem, ale pełny model korzysta również z informacji zastępczej zawartej w pozostałych rankingach. Maskowanie cech w danych treningowych nie poprawiło wyniku, dlatego nie zostało włączone do finalnego estymatora. Ablacja pełni jednak ważną funkcję kontrolną: pokazuje, że model nie jest arbitralną czarną skrzynką, lecz wykorzystuje przewidywalną strukturę sygnałów rankingowych.

Osobno sprawdzono, czy do finalnego wejścia warto dodać predykcje i parametry rankingów drużynowych. Wariant rozszerzony Player+Team W20-Binomial uzyskał minimalnie niższy LogLoss niż wariant player-based, ale różnica była bardzo mała i nie została potwierdzona w miesięcznym bootstrapie blokowym. Wynik przedstawiono w tabeli~\ref{tab:player_team_ablation}. Z tego powodu finalny model zachowuje prostszą strukturę: komponent rankingowy opiera się na sygnałach zawodniczych, natomiast informacja drużynowa jest reprezentowana przez pełny kontekst historyczny W20.

\begin{table}[H]
    \centering
    \caption{Ablacja dodania rankingów drużynowych do finalnego wejścia W20-Binomial}
    \label{tab:player_team_ablation}
    \begin{tabular}{lrrrr}
        \toprule
        Wariant & AUC & LogLoss & Brier & ECE \\
        \midrule
        Player-based W20-Binomial & 0{,}7553 & 0{,}5858 & 0{,}2007 & 0{,}0131 \\
        Player+Team W20-Binomial & 0{,}7554 & 0{,}5857 & 0{,}2007 & 0{,}0144 \\
        \bottomrule
    \end{tabular}
\end{table}

W miesięcznym bootstrapie blokowym różnica $\Delta=\mathrm{LogLoss}_{\mathrm{Player+Team}}-\mathrm{LogLoss}_{\mathrm{Player}}$ wyniosła $-0{,}0001$, a przedział ufności 95\% był równy $[-0{,}0002,\ 0{,}0020]$. Ponieważ przedział obejmuje zero, dodatkowe cechy team-based potraktowano jako zwiększenie złożoności bez stabilnej poprawy jakości predykcji.

\begin{figure}[H]
    \centering
    \includegraphics[width=0.9\textwidth]{logistic_masking_ablation_bootstrap/masking_monthly_block_bootstrap_ci.png}
    \caption{Miesięczny bootstrap blokowy różnic LogLoss wariantów maskingu względem pełnego modelu kontekstowego.}
    \label{fig:masking_block_bootstrap}
\end{figure}

\section{Symetryzacja kolejności drużyn}

\subsection{Źródło asymetrii technicznej}

Regresja logistyczna trenowana na wektorze cech zależnym od kolejności stron meczu generuje predykcję dla zdarzenia, że wygra strona zapisana jako Team 1. Problem ten jest technicznie naturalny, ponieważ dane wejściowe mają postać różnic lub par cech Team 1--Team 2. Jednocześnie z perspektywy sportowej prawdopodobieństwo powinno spełniać warunek antysymetrii: jeżeli ten sam mecz zostanie zapisany w odwróconej kolejności, to prawdopodobieństwo zwycięstwa pierwszej strony powinno być dopełnieniem poprzedniej predykcji.

\subsection{Procedura uśredniania predykcji}

W praktyce sprawdzono, że surowy estymator nie jest idealnie niewrażliwy na kolejność drużyn. Dla każdego meczu wygenerowano więc dwie predykcje: predykcję w oryginalnej orientacji oraz predykcję po zamianie stron. Predykcję po zamianie stron przeliczono z powrotem na prawdopodobieństwo zwycięstwa pierwotnej strony, a następnie uśredniono oba oszacowania:

\begin{equation}
    p_{sym}(A>B) = \frac{p_{orig}(A>B) + \left(1 - p_{swap}(B>A)\right)}{2}.
    \label{eq:symmetrized_prediction}
\end{equation}

Tak zdefiniowany wariant zachowuje informację wyuczoną przez regresję logistyczną, ale ogranicza artefakt wynikający z arbitralnego zapisu meczu. Diagnostyka wrażliwości kolejności wykazała średnią bezwzględną różnicę prawdopodobieństw równą około 0{,}0287, medianę 0{,}0237 oraz maksimum 0{,}2113. Po symetryzacji surowy model uzyskał AUC 0{,}7561 i LogLoss 0{,}5858, nieznacznie poprawiając rozróżnianie względem predykcji w oryginalnej orientacji. Wniosek z tej analizy jest metodologiczny: model końcowy powinien respektować symetrię problemu, nawet jeżeli sama poprawa metryk jest niewielka.

\section{Kalibracja prawdopodobieństw}

\subsection{Schemat expanding-window}

Ponieważ celem pracy jest predykcja probabilistyczna, końcowy wariant modelu poddano dodatkowej analizie kalibracji. Porównano predykcje surowe, kalibrację Platta oraz kalibrację izotoniczną. Obie metody były stosowane w schemacie expanding-window: kalibrator dla danego folda uczono wyłącznie na wcześniejszych foldach, bez wykorzystania przyszłych obserwacji. Dzięki temu kalibracja zachowuje chronologiczny charakter ewaluacji.

\subsection{Porównanie Platta i kalibracji izotonicznej}

Kalibracja Platta polegała na dopasowaniu regresji logistycznej do logitu surowych prawdopodobieństw modelu. Jej zaletą jest mała liczba parametrów i odporność na mniejsze próbki kalibracyjne. Kalibracja izotoniczna jest bardziej elastyczna, ale może łatwiej dopasować lokalny szum, szczególnie przy dryfie danych i nierównej liczbie obserwacji w kolejnych okresach.

\begin{table}[H]
    \centering
    \caption{Porównanie wariantów symetryzacji i kalibracji na wspólnej próbie rynkowej}
    \label{tab:symmetry_calibration_comparison}
    \small
    \begin{tabular}{lrrrr}
        \toprule
        Wariant & AUC & LogLoss & Brier & ECE \\
        \midrule
        Symetryczny + Platt expanding & 0{,}7551 & 0{,}5854 & 0{,}2007 & 0{,}0100 \\
        Oryginalny + Platt expanding & 0{,}7542 & 0{,}5862 & 0{,}2010 & 0{,}0098 \\
        Symetryczny surowy & 0{,}7551 & 0{,}5867 & 0{,}2011 & 0{,}0193 \\
        Oryginalny surowy & 0{,}7543 & 0{,}5868 & 0{,}2011 & 0{,}0125 \\
        Symetryczny + isotonic expanding & 0{,}7532 & 0{,}5901 & 0{,}2013 & 0{,}0097 \\
        Oryginalny + isotonic expanding & 0{,}7519 & 0{,}5906 & 0{,}2018 & 0{,}0104 \\
        \bottomrule
    \end{tabular}
\end{table}

\begin{figure}[H]
    \centering
    \includegraphics[width=0.92\textwidth]{final_symmetric_calibrated_market_comparison/final_symmetric_calibrated_calibration_logloss.png}
    \caption{Porównanie wariantów symetryzacji i kalibracji według LogLoss na wspólnej próbie rynkowej.}
    \label{fig:final_calibration_logloss}
\end{figure}

\begin{figure}[H]
    \centering
    \includegraphics[width=0.92\textwidth]{final_symmetric_calibrated_market_comparison/final_symmetric_calibrated_calibration_ece.png}
    \caption{Porównanie wariantów symetryzacji i kalibracji według ECE na wspólnej próbie rynkowej.}
    \label{fig:final_calibration_ece}
\end{figure}

Najniższy LogLoss uzyskał wariant symetryczny z kalibracją Platta. Kalibracja izotoniczna dawała bardzo niskie ECE, ale pogarszała LogLoss i AUC. Zgodnie z hierarchią metryk przyjętą w pracy wybrano więc wariant Platta: poprawiał kalibrację względem surowego modelu, a jednocześnie nie obniżał jakości probabilistycznej mierzonej główną metryką. Ostateczny model kierowany do porównania z benchmarkami rynkowymi to \textbf{skalibrowany symetryczny LR-ElasticNet-W20-Binomial}.

\section{Wnioski z etapu metamodelowania}

Etap metamodelowania prowadzi do sześciu głównych wniosków. Po pierwsze, cechy zawodnicze pozostają najważniejszą podstawą predykcji, a Player Glicko-2 jest najsilniejszym pojedynczym sygnałem rankingowym. Po drugie, pełny kontekst W20 uzupełnia rankingi o informacje o aktualnej formie, stylu, ekonomii, obiektach i tempie gry. Po trzecie, format serii Bo1/Bo3/Bo5 najlepiej uwzględnić strukturalnie przez transformację dwumianową prawdopodobieństw rankingowych. Po czwarte, mimo przetestowania modeli drzewiastych i MLP, najlepiej generalizującym estymatorem okazała się regularyzowana regresja logistyczna ElasticNet. Po piąte, finalny model powinien respektować symetrię problemu względem kolejności drużyn. Po szóste, kalibracja Platta poprawia jakość prawdopodobieństw bez pogorszenia LogLoss, dlatego została użyta w wariancie końcowym.

Najważniejszym rezultatem tego rozdziału jest wybór modelu głównego do porównania z benchmarkami bukmacherskimi: \textbf{skalibrowanego symetrycznego LR-ElasticNet-W20-Binomial}. Model ten będzie w następnym rozdziale oceniony na wspólnej próbie meczów posiadających zarówno dane sportowe, jak i kursy bukmacherskie.
